\begin{textsample}{POS Dim 2 – al – Score 23.00 – al/al\_em\_9.txt}  \label{ex:f2_pos_254}
POSSIBILIDADES DE ARQUITETURA CURRICULAR PARA O ENSINO MÉDIO A motivação para a construção de uma proposta de (re)desenho curricular para as redes de ensino, foi a necessidade de implementação de um Ensino Médio que despertasse maior interesse, de forma a fazer sentido na vida dos estudantes.
A argumentação que ampara todas as discussões para a construção deste documento é a Base Nacional Comum Curricular e todas as \textbf{diretrizes} que foram alicerces da (re)elaboração dos documentos curriculares dos entes federados.
A necessidade do entendimento de como fazer um \textbf{currículo} estimulante e que prepare jovens, mobilizou uma frente de trabalho, composta por estudiosos e representantes de todos os estados e do Distrito Federal, para construir possibilidades de arquitetura curricular.
A Lei 13.415 de 2017 destaca, em seu Art. 35 -A, § 7º, que : > “ Os \textbf{currículos} do ensino médio deverão considerar a formação integral do aluno, de maneira a adotar um trabalho voltado para a construção de seu projeto de vida e para sua formação nos aspectos físicos, cognitivos e socioemocionais. ” Nessa perspectiva, é necessário um ensino médio com uma arquitetura curricular \textbf{flexível}, que promova o protagonismo juvenil, a organização interdisciplinar entre os componentes curriculares, proporcionando assim uma visão de integração e complementaridade, onde o estudante possa realizar seu projeto de vida.
A partir do marco legal para o novo Ensino Médio, os \textbf{currículos} passam a ser compostos pela formação geral básica e pelos \textbf{itinerários} formativos de maneira indissociável, como preconizado nas DCNEM/2018 ( \textbf{Parecer} CNE/CEB nº 3/2018 e Resolução CNE/CEB nº 3/2018 ).
A Formação Geral Básica deve ter, no máximo, 1800 horas, e os \textbf{Itinerários} Formativos devem ser desenvolvidos em, no mínimo, 1200 horas anuais.
Compõem os \textbf{Itinerários} Formativos o Projeto de Vida, os Aprofundamentos e as \textbf{Eletivas}.
Arquitetura do Ensino Médio A Formação Geral Básica é definida, nos termos das \textbf{Diretrizes} Curriculares Nacionais para o Ensino Médio, em seu Art. 6º, Inciso II, como : conjunto de competências e habilidades das áreas de conhecimento \textbf{previstas} na Base Nacional Comum Curricular ( BNCC ), que aprofundam e consolidam as aprendizagens essenciais do ensino fundamental, a compreensão de problemas complexos e a reflexão sobre soluções para eles ; (... ) é composta por competências e habilidades \textbf{previstas} na Base Nacional Comum Curricular ( BNCC ) e articuladas como um todo indissociável, enriquecidas pelo contexto histórico, econômico, social, ambiental, cultural local, do mundo do trabalho e da prática social, e deverá ser organizada por áreas de conhecimento : I - linguagens e suas tecnologias ; II - matemática e suas tecnologias ; III - ciências da natureza e suas tecnologias ; IV - ciências humanas e sociais aplicadas. (... ) O \textbf{itinerário} formativo é definido como um conjunto de situações e atividades educativas que os estudantes podem fazer a partir das escolhas conforme seus interesses, cujo objetivo central é aprofundar e/ou ampliar aprendizagens em uma ou mais áreas do conhecimento e/ou na Formação \textbf{técnica} e Profissional ( Referenciais Curriculares para elaboração de \textbf{itinerários} formativos : \textbf{Portaria} nº 1.432/2018 ).
O termo \textbf{itinerário}, cuja etimologia expressa a origem no latim itinerarius.a.um, faz a ligação sobre “ percurso que se pretende seguir... ”, “ parte do caminho percorrido... ”.
E pelo termo, não podemos deixar de ressaltar que o caminho percorrido e a ser percorrido pelo estudante na sua trajetória de vida está dentro de uma caminhada chamada Educação Básica.
Assim, não podemos dissociar o Ensino Fundamental do Ensino Médio, pois um dará subsídios ao outro dentro de uma formação cidadã.
Os \textbf{itinerários} formativos devem considerar as demandas e necessidades do mundo contemporâneo, sempre considerar o contexto local e estar em sintonia com os interesses dos estudantes e de forma a promover sua inserção na sociedade.
As redes podem \textbf{ofertar} \textbf{Itinerários} Formativos simples, estruturados com foco em uma área do conhecimento ou na formação \textbf{técnica} e profissional, ou \textbf{Itinerários} Formativos integrados, que mobilizam competências e habilidades de uma ou mais áreas do conhecimento ou destas com a formação \textbf{técnica} e profissional.
A Resolução CNE/CP nº 1, de 5 de \textbf{Janeiro} de 2021, em seu Art. 5º, § 5º, define o \textbf{itinerário} formativo na Educação Profissional e Tecnológica como o conjunto de unidades curriculares, etapas ou módulos que compõem a sua organização em eixos tecnológicos e respectiva área tecnológica ( \textbf{Diretrizes} Curriculares Nacionais Gerais para a Educação Profissional e Tecnológica, 2021 ).
O \textbf{itinerário} de Formação \textbf{Técnica} e Profissional, associado às demais áreas do conhecimento ou não, propiciará ao estudante ser protagonista do seu \textbf{currículo} formativo, podendo escolher trilhas formativas, cujo objetivo é alinhar os seus conhecimentos da formação geral básica, buscando ajustar ao seu projeto de vida e interesses.
Entretanto, esse \textbf{itinerário} não encerra em si a formação cidadã, visto que os interesses dos estudantes devem ser despertados para transcender a formação \textbf{técnica} e buscar o mundo do trabalho como uma forma de significar a sua vida para além da sua sustentabilidade pessoal.
Sendo assim, a construção deste \textbf{itinerário} não deverá ser pensada para o trabalho tão somente, mas para que cada estudante possa relacionar -se com o mundo e seus significados.
De acordo com as \textbf{Diretrizes} Curriculares Nacionais para o Ensino Médio, os \textbf{itinerários} precisam ser organizados de forma a contemplar um ou os quatro Eixos Estruturantes, a saber : Investigação Científica, Processos Criativos, Mediação e Intervenção Sociocultural e Empreendedorismo.
No entanto, considerando que os quatro eixos estruturantes são complementares, recomenda -se que os \textbf{Itinerários} Formativos incorporem e integrem todos eles, para proporcionar aos estudantes experienciar diversas situações que resultem em aprendizagens relevantes para sua formação integral.
Para que os \textbf{itinerários} formativos possam alcançar seus objetivos, fazendo sentido nas escolhas dos estudantes, é necessário organizar a \textbf{oferta} de modo a : aprofundar as aprendizagens relacionadas às competências gerais das áreas do conhecimento ; consolidar a formação integral proporcionando a autonomia para realização do projeto de vida dos estudantes ; promover a incorporação de valores como ética, liberdade, democracia, justiça social, pluralidade, solidariedade e sustentabilidade ; promover atividades que desenvolvam as habilidades que permitam aos estudantes uma visão de mundo ampla e diversa, onde possam saber tomar decisões e agir nas mais diversas situações em seus territórios de vivência social.
As possibilidades de organização curricular, incluindo a \textbf{oferta} de \textbf{itinerários} formativos favorecem o uso de metodologias e estratégias didáticas que promovam a interdisciplinaridade e que possibilitem aos estudantes a análise de situações-problema no seu cotidiano, reflexão crítica e tomada de decisão responsável.
Para a distribuição da \textbf{carga} \textbf{horária}, as redes devem considerar suas especificidades.
Diversos arranjos são possíveis, desde que seja \textbf{ofertada} a \textbf{carga} \textbf{horária} mínima \textbf{prevista}, \textbf{atendendo} à formação geral básica, com máximo de 1.800 horas, e o \textbf{itinerário} formativo, com mínimo de 1.200 horas, totalizando assim 3.000 horas de formação.
Assim, os exemplos abaixo demonstram algumas das possíveis maneiras de distribuição das 3.000 horas da \textbf{carga} \textbf{horária} do Ensino Médio : Dentre os diversos arranjos possíveis, apresentamos como sugestão para a implementação uma distribuição da \textbf{carga} \textbf{horária} em que, no 1º ano do ensino médio, as unidades curriculares referentes à formação geral básica somam 800 horas anuais, enquanto que o \textbf{itinerário} formativo, a parte \textbf{flexível} do \textbf{currículo}, começa a ser abordado em 200 horas anuais, com a construção do Projeto de Vida do estudante e algumas \textbf{eletivas}, perfazendo uma \textbf{carga} anual mínima de 1.000 horas.
Nos anos seguintes, à medida que decresce a \textbf{carga} \textbf{horária} \textbf{destinada} à Formação Geral Básica aumenta aquela \textbf{destinada} ao \textbf{itinerário} formativo.
A partir do 2º ano, o estudante do ensino médio \textbf{cursa}, além das \textbf{eletivas} e Projeto de Vida, as unidades curriculares que promovem o aprofundamento ou ampliação dos conhecimentos da área escolhida.
Nesse arranjo, o ingresso do estudante no aprofundamento ocorre após um processo de orientação e apoio à construção de seu projeto de vida.
Ensino Médio - Proposta para o Sistema Estadual de Ensino de Alagoas O Projeto de Vida é um trabalho pedagógico intencional e estruturado para desenvolver a capacidade do estudante de dar sentido à sua existência, tomar decisões, planejar o futuro e agir no presente com autonomia e responsabilidade.
Deve promover aos estudantes o autoconhecimento, a identificação de seus interesses, desejos e aspirações, o reconhecimento e desenvolvimento de suas potencialidades, definindo seus objetivos e estabelecendo estratégias e metas para alcançá -los.
O Projeto de Vida permite aos estudantes construírem suas trajetórias para que atinjam a realização pessoal, profissional e contribuam positivamente com o meio em que vivem.
Nessa perspectiva, não deve ser confundido apenas com escolha profissional, pois está intrinsecamente relacionado às perspectivas futuras dos jovens em âmbitos diversos.
O Projeto de Vida deve estar presente nos três anos do Ensino Médio com uma unidade curricular.
É um processo educativo que deve estar alinhado com a comunidade escolar e ocorrer, além da sala de aula, nos diversos momentos na escola, de forma a promover o desenvolvimento de competências como autoconfiança, determinação e resiliência, dentre outros.
Por seu caráter eminentemente transversal à trajetória do estudante no ensino médio, bem como à trajetória já construída na educação básica, o trabalho com o Projeto de Vida deve estar sempre articulado e integrado com todas as propostas de ensino presentes no \textbf{itinerário} formativo e formação geral.
O Referencial Curricular de Alagoas tem como um de seus objetivos principais, apontar possibilidades de execução de arquiteturas curriculares que possam ser adequadas às mais diferentes realidades do território, buscando levar os mais variados conhecimentos que sejam capazes de correlacionar a vida dos jovens ao futuro tão almejado e planejado.
Um futuro onde o jovem possa se ver como um membro da sociedade atuando como um protagonista na transformação de uma realidade que, por vezes, é vista de forma muito dura, seca e sem perspectivas.
Para um melhor entendimento, apresentamos o quadro a seguir, que detalha distribuição de \textbf{carga} \textbf{horária}, conforme a implementação da proposta de \textbf{oferta} de ensino médio.
Enquanto estratégia de mobilizar as \textbf{ofertas} e propostas de Ensino Médio no sistema estadual de ensino em Alagoas, recomendamos um modelo que apresenta \textbf{carga} \textbf{horária} de 80h anual, pelo menos, para todas unidades curriculares.
Neste modelo nenhuma unidade terá menos de 2 horas semanais.
Além de considerar que, a cada ano do Ensino Médio, se disponha \textbf{carga} \textbf{horária} para unidades curriculares \textbf{Eletivas}, \textbf{obrigatoriamente}.
Também, apresentamos a unidade curricular Projeto de Vida, obrigatória, e em todos os anos do ensino médio.
O \textbf{itinerário} formativo apresenta o primeiro ano sem aprofundamento, com a perspectiva de intensificar o trabalho da BNCC, todavia o segundo e terceiro anos avançam progressivamente no aprofundamento deste IF.
Ressalta -se a importância da SERIAÇÃO ou distribuição das unidades curriculares da BNCC ou Formação Geral Básica.
Desta forma, é importante definir quais unidades curriculares vão estar em cada ano e como serão articuladas para o desenvolvimento da BNCC e IF.
No modelo recomendado, o tempo de desenvolvimento da BNCC é inverso ao tempo de desenvolvimento do Itinerário Formativo, isto é, a medida que a \textbf{carga} \textbf{horária} do Itinerário Formativo é progressivo, a da BNCC é regressiva.

% matched lemmas: atender, carga, currículo, cursar, destinar, diretriz, eletivo, flexível, horário, itinerário, janeiro, obrigatoriamente, oferta, ofertar, parecer, portaria, prever, técnico
\end{textsample}
